\chapter{Results and Discussion}
\label{chap:results}

% This chapter is intentionally a results-ready template. Replace every
% \textit{Pending official run} cell only from 05_metric_summary.json and the
% validated 03_judged_traces.json produced by the official 100-query run.

\section{Overview}

This chapter reports the performance of the final v2.5 Personalized Travel
Guide Assistant on the official 100-query multi-turn dataset. Results are
organized by query understanding, retrieval, evidence recovery, final-answer
quality, and operational efficiency. Stage-level findings are interpreted
together because a final answer may be limited by the corpus, parsing,
retrieval, coverage checking, or generation.

At the time this chapter draft was prepared, the official run had not yet been
finalized and human-validated. Consequently, the tables contain explicit
placeholders rather than smoke-test values. Development results must not be
reported as official findings.

\section{Dataset and Run Integrity}

Table~\ref{tab:run-integrity} summarizes the completed evaluation artifacts.

\begin{table}[H]
\centering
\caption{Official run integrity summary}
\label{tab:run-integrity}
\begin{tabular}{ll}
\hline
\textbf{Field} & \textbf{Value} \\
\hline
Expected cases & 100 \\
Pipeline-complete cases & \textit{Pending official run} \\
Judge-complete cases & \textit{Pending official run} \\
Human-reviewed cases & \textit{Pending human review} \\
Pipeline version & \texttt{personalized-rag-eval-v2.5} \\
Corpus fingerprint & \textit{Pending official run} \\
Parser fallback cases & \textit{Pending official run} \\
\hline
\end{tabular}
\end{table}

Before interpreting aggregate metrics, all expected case identifiers should be
present exactly once, and no case should contain an incompatible judge schema,
missing relevance grades, or missing final-answer scores.

\section{Query-Understanding Results}

Table~\ref{tab:understanding-results} reports the query-understanding metrics.

\begin{table}[H]
\centering
\caption{Query-understanding performance}
\label{tab:understanding-results}
\begin{tabular}{lc}
\hline
\textbf{Metric} & \textbf{Score} \\
\hline
Intent Accuracy & \textit{Pending official run} \\
Operation Accuracy & \textit{Pending official run} \\
Query Constraint F1 & \textit{Pending official run} \\
Retrieval Facet F1 & \textit{Pending official run} \\
\hline
\end{tabular}
\end{table}

Intent and operation should be analyzed separately. A correct operation with an
incorrect intent may still produce a superficially appropriate plan but apply
the wrong category assumptions. Constraint and facet errors should be divided
into omissions, unsupported additions, vocabulary-normalization disagreements,
and geographic type errors.

The discussion should report the most frequent confusion pairs and state
whether each disagreement reflects a parser error or an ambiguous synthetic
reference label. Particular attention should be given to recommendation versus
accommodation-search requests, broad region names, negative constraints, and
follow-up queries whose meaning depends on conversation state.

\section{Retrieval Results}

\begin{table}[H]
\centering
\caption{Retrieval ranking performance at depth five}
\label{tab:retrieval-results}
\begin{tabular}{lc}
\hline
\textbf{Metric} & \textbf{Score} \\
\hline
nDCG@5 & \textit{Pending official run} \\
Precision@5 & \textit{Pending official run} \\
Mean judged relevance grade & \textit{Pending official run} \\
Queries with no grade-2/3 chunk in top five & \textit{Pending official run} \\
\hline
\end{tabular}
\end{table}

nDCG@5 measures whether the most useful evidence appears near the top, while
Precision@5 measures how much of the evaluated context is useful. Cases with
high Precision@5 but lower nDCG@5 indicate that useful chunks were retrieved but
ordered poorly. Low values for both metrics indicate either a retrieval failure,
overly restrictive filtering, query decomposition error, or missing corpus
content.

Results should also be grouped by query type and complexity. Recommended groups
include single-destination lookup, recommendation, comparison, itinerary,
transport, accommodation, food, and multi-turn follow-up. This analysis can
show whether the global average hides systematic weaknesses in multi-aspect or
multi-destination requests.

\section{Evidence Coverage and Recovery}

\begin{table}[H]
\centering
\caption{Coverage, readiness, and recovery behavior}
\label{tab:coverage-results}
\begin{tabular}{lc}
\hline
\textbf{Measure} & \textbf{Value} \\
\hline
Mean coverage ratio before recovery & \textit{Pending official run} \\
Mean coverage ratio after recovery & \textit{Pending official run} \\
Cases with recovery performed & \textit{Pending official run} \\
Cases whose coverage improved & \textit{Pending official run} \\
Mean new candidates per recovery case & \textit{Pending official run} \\
Complete readiness & \textit{Pending official run} \\
Partial readiness & \textit{Pending official run} \\
Insufficient readiness & \textit{Pending official run} \\
\hline
\end{tabular}
\end{table}

The before/after comparison measures whether targeted retrieval contributes new
evidence within the same final pipeline. It does not constitute a separate
legacy ablation. Recovery should be considered successful only when it adds
useful selected evidence or improves requirement coverage, not merely when it
returns additional chunks.

Unresolved requirements should be grouped by diagnostic label:
\texttt{not\_probed}, \texttt{likely\_corpus\_gap},
\texttt{likely\_corpus\_content\_gap},
\texttt{likely\_reranking\_failure}, and
\texttt{selected\_evidence\_insufficient}. A large number of
\texttt{not\_probed} cases indicates that the recovery limit or checker query
selection requires attention; it does not prove that the corpus is incomplete.

\section{Final-Answer Results}

\begin{table}[H]
\centering
\caption{Mean final-answer scores on the five-point scale}
\label{tab:answer-results}
\begin{tabular}{lc}
\hline
\textbf{Metric} & \textbf{Mean score} \\
\hline
Correctness & \textit{Pending official run} \\
Faithfulness & \textit{Pending official run} \\
Personalization Adherence & \textit{Pending official run} \\
Completeness & \textit{Pending official run} \\
\hline
\end{tabular}
\end{table}

The final scores should be interpreted with readiness. For example, a partial
answer may correctly acknowledge a missing requirement and therefore remain
faithful, while receiving a lower completeness score. Conversely, a complete
and detailed answer with unsupported claims may obtain high apparent coverage
but low faithfulness.

Personalization errors should be divided into missed relevant preferences,
incorrectly applied irrelevant preferences, violations of avoidance
constraints, and style mismatches. Safety-sensitive dietary cases require
manual inspection because an unsupported claim of safety is more serious than
a general recommendation mismatch.

\section{Relationship Between Retrieval and Answer Quality}

To answer the central evaluation question, cases should be grouped or plotted
by retrieval and answer scores. The following relationships are particularly
useful:

\begin{itemize}
    \item nDCG@5 versus correctness;
    \item Precision@5 versus faithfulness;
    \item coverage ratio versus completeness;
    \item readiness mode versus all four answer scores; and
    \item recovery improvement versus change in supported requirements.
\end{itemize}

These comparisons should not be interpreted as proof of causality. They help
determine whether answer limitations are associated with evidence quality. A
low answer score despite strong retrieval suggests generation or personalization
failure, whereas weak retrieval and weak completeness suggest an upstream
evidence problem.

\section{Latency, Token Usage, and Estimated Cost}

\begin{table}[H]
\centering
\caption{Operational performance of the official run}
\label{tab:efficiency-results}
\begin{tabular}{lc}
\hline
\textbf{Measure} & \textbf{Value} \\
\hline
Mean pipeline latency per case & \textit{Pending official run} \\
Median pipeline latency per case & \textit{Pending official run} \\
95th-percentile pipeline latency & \textit{Pending official run} \\
Total pipeline tokens & \textit{Pending official run} \\
Total judge tokens & \textit{Pending official run} \\
Estimated generation cost & \textit{Pending official run} \\
Estimated pipeline cost & \textit{Pending official run} \\
Estimated judging cost & \textit{Pending official run} \\
Estimated total experiment cost & \textit{Pending official run} \\
\hline
\end{tabular}
\end{table}

Stage-level latency should identify the dominant operations, especially query
parsing, planning, coverage checking, recovery, answer generation, embedding,
and reranking. Estimated cost must be reported together with the pricing
snapshot and described as an estimate rather than an invoice.

\section{Error Analysis}

Representative failures should be selected from the trace using a consistent
procedure rather than anecdotal preference. Recommended categories are:

\begin{enumerate}
    \item parser omission or over-extraction;
    \item geographic normalization or filtering error;
    \item planner under-decomposition or over-decomposition;
    \item relevant evidence retrieved but reranked too low;
    \item unbalanced evidence across tasks or destinations;
    \item likely corpus gap;
    \item coverage-checker disagreement;
    \item unsupported answer claim;
    \item missed or excessive personalization; and
    \item malformed model output handled by fallback.
\end{enumerate}

For each selected case, the report should include the query, relevant profile
or state, parsed representation, planned requirements, top evidence, readiness,
answer excerpt, scores, diagnosed cause, and proposed improvement. Personally
identifiable information and API metadata must be removed.

\section{Discussion by Research Question}

\subsection{RQ1: Memory and Conversation Context}

Insert an evidence-based answer using cross-conversation personalization cases,
follow-up rewriting, state preservation, and errors involving irrelevant or
temporary preferences. State clearly whether the evidence supports effective
separation of persistent and temporary context.

\subsection{RQ2: Hybrid Retrieval and Reranking}

Use nDCG@5, Precision@5, query-type breakdowns, filter relaxations, and
reranking failures to explain how well the combined retrieval strategy works.
Avoid claiming that hybrid retrieval is superior to every alternative unless a
controlled comparison was actually conducted.

\subsection{RQ3: Coverage and Recovery}

Use coverage before and after recovery, newly covered requirements, readiness,
and diagnostic categories. Distinguish successful recovery from cases where
the corpus does not provide the requested evidence.

\subsection{RQ4: Multi-stage Evaluation}

Discuss whether stage-level traces explained final-answer outcomes that would
have been ambiguous under final-answer scoring alone. Include the operational
cost of the additional planner, checker, and judge calls.

\section{Limitations and Validity of the Findings}

The reported findings are conditioned on the official corpus snapshot,
synthetic dataset, model versions, prompts, and judge. Synthetic queries provide
control but cannot represent all real travelers. The independent judge reduces
self-evaluation but remains an LLM with known biases. Human review improves
confidence but does not make the relevance pool exhaustive; therefore Recall@5
is not reported.

The evaluation covers only the final proposed system and internal recovery
behavior. It should not make causal claims about the superiority of every
component over an unimplemented baseline. Broader user studies and controlled
component comparisons remain future work.

\section{Chapter Summary}

Replace this paragraph after the official run with a concise summary containing
the principal understanding, retrieval, final-answer, recovery, and efficiency
findings. Report only values that can be reproduced from the exported metric
summary and validated judged traces.

